\section*{الفصل \ref{ch:g4:symmetry} --- التماثل}

\begin{solution}{exo:g4:symmetry:1}
المربع: نعم ($4$ محاور). والمستطيل: نعم ($2$). والمثلث المختلف
الأضلاع: لا طيّ ينجح فيه. والدائرة: نعم --- أي مستقيم يمرّ بالمركز.
\end{solution}

\begin{solution}{exo:g4:symmetry:2}
المربع: $4$ (خطّا المنتصف، والقطران). والمستطيل: $2$ (خطّا
المنتصف لا غير). والمثلث المتساوي الأضلاع: $3$ (واحد يمرّ بكل
رأس).
\end{solution}

\begin{solution}{exo:g4:symmetry:3}
بالطيّ على \omterm{def:g4:geometry:circle}{القطر}، يكون للنصفين المثلثين الشكل
نفسه لكنهما لا ينطبقان --- فأحدهما «مقلوب الاتجاه»:
فينزل النقل بجانب الشكل لا عليه. إذن \omterm{def:g4:geometry:circle}{القطر} ليس
محورًا.
\end{solution}

\begin{solution}{exo:g4:symmetry:4}
صورة الحرف L تقع على بُعد $2$ من المربعات من الجهة
\emph{الأخرى} للمحور، مكتوبة معكوسة (حرف L مرآتي)، بالأحجام نفسها.
\end{solution}

\begin{solution}{exo:g4:symmetry:5}
يُظهر الشكل المكتمل العلم وانعكاسه مقلوبًا
تحت المحور، السارية تحت السارية.
\end{solution}

\begin{solution}{exo:g4:symmetry:6}
لا تتحرك $P$: فالطيّ على المحور يترك كل نقطة
\emph{على} المحور في مكانها --- فالنقطة $P$ صورة نفسها.
\end{solution}

\begin{solution}{exo:g4:symmetry:7}
المحور الرأسي: A و H و O و T. والمحور الأفقي: B و E و H و O.
وكلاهما: H و O. ولا محور: F و N.
\end{solution}

\begin{solution}{exo:g4:symmetry:8}
نعم: فالمثلث متساوي الساقين، ومحوره يمرّ بالرأس
الواقع بين ضلعَي $5$ سم وبمنتصف القاعدة $3$ سم.
ويبدّل الطيّ بين الضلعين المتساويين.
\end{solution}

\begin{solution}{exo:g4:symmetry:9}
كل شكل له محوران بالضبط يصلح: مستطيل (غير مربع)، أو
شكل زائد ذراعاه بطولين مختلفين (أفقي طويل، ورأسي
قصير)، أو بيضاوي. والمحوران يتقاطعان دائمًا في مركز
الشكل.
\end{solution}

\begin{solution}{exo:g4:symmetry:10}
البيت الكامل هو النصف الأيسر مع توأمه المرآتي: جدار
عرضه ضعف نصف الجدار، تحت سقف مثلث كامل. وكل
طول على اليمين يساوي شريكه على اليسار --- فالتماثل
ينسخ الأطوال بالضبط.
\end{solution}

\begin{solution}{exo:g4:symmetry:11}
الرسوم تؤيد الاتجاهين: فالمثلث المتساوي الساقين ينطوي
على المستقيم المارّ برأسه (ضلعان متساويان $\to$ محور)؛ والمتساوي
الأضلاع ينطوي بثلاث طرق؛ والمختلف الأضلاع لا محور له (لا
أضلاع متساوية $\to$ لا محور). وادّعاء آية المزدوج صحيح --- وسيُبرهن
عليه في \cref{ch:g6:symmetry}.
\end{solution}
